\documentclass[12pt, letterpaper]{article}

% ── Paquetes ────────────────────────────────────────────────────────────────
\usepackage[utf8]{inputenc}
\usepackage[T1]{fontenc}
\usepackage[spanish]{babel}
\usepackage[margin=2.5cm]{geometry}
\usepackage{lmodern}
\usepackage{microtype}
\usepackage{parskip}
\usepackage{titlesec}
\usepackage{titling}
\usepackage{hyperref}
\usepackage{xcolor}
\usepackage{enumitem}
\usepackage{booktabs}
\usepackage{array}
\usepackage{fancyhdr}
\usepackage{graphicx}
\usepackage{amsmath}
\usepackage[most]{tcolorbox}
\usepackage{xcolor}
\usepackage{parskip}
\usepackage{datetime}

% ── Colores ─────────────────────────────────────────────────────────────────
\definecolor{primary}{HTML}{1B4F72}
\definecolor{accent}{HTML}{2E86C1}
\definecolor{lightgray}{HTML}{F2F3F4}
\definecolor{darkgray}{HTML}{555555}

% ── Hipervínculos ────────────────────────────────────────────────────────────
\hypersetup{
    colorlinks=true,
    linkcolor=accent,
    urlcolor=accent,
    citecolor=accent,
    pdftitle={Propuesta de Proyecto: Prediusly},
    pdfauthor={Equipo Prediusly},
}

% ── Encabezado y pie de página ───────────────────────────────────────────────
\pagestyle{fancy}
\fancyhf{}
\fancyhead[L]{\small\color{darkgray}\textit{Prediusly — Predicción de Precios de Viviendas y Predios}}
\fancyhead[R]{\small\color{darkgray}\thepage}
\fancyfoot[C]{\small\color{darkgray}Propuesta de Proyecto · \the\year}
\renewcommand{\headrulewidth}{0.4pt}

% ── Formato de secciones ─────────────────────────────────────────────────────
\titleformat{\section}
  {\large\bfseries\color{primary}}
  {\thesection.}{0.6em}{}[\vspace{-0.3em}\rule{\linewidth}{0.6pt}]

\titleformat{\subsection}
  {\normalsize\bfseries\color{accent}}
  {\thesubsection.}{0.5em}{}

\titlespacing{\section}{0pt}{1.4em}{0.6em}
\titlespacing{\subsection}{0pt}{0.9em}{0.3em}

% ── Título ───────────────────────────────────────────────────────────────────
\title{%
  {\LARGE\bfseries\color{primary} Prediusly}\\[0.3em]
  {\large\color{darkgray} Predicción de Precios de Viviendas y Predios\\
   en Áreas Metropolitanas mediante Datos Geoespaciales}\\[0.8em]
  \rule{\linewidth}{1.2pt}\\[0.4em]
  {\normalsize\color{darkgray} Propuesta de Proyecto}
}
\author{\color{darkgray} Equipo Prediusly}
\date{\color{darkgray}10 de abril de 2026}

% ════════════════════════════════════════════════════════════════════════════
\begin{document}

\maketitle
\thispagestyle{empty}

\vspace{0.5em}
{\centering\color{darkgray}
\textbf{Equipo Prediusly}\\
Cristobal Robinson Morales | Nicole Russo Martínez \\
Giovanni Jimenez Andrade | Bastián Medina Gómez\par
}

\vspace{1em}


% ════════════════════════════════════════════════════════════════════════════
\section{Contexto y Motivación}

El precio de las viviendas y predios urbanos es el resultado de una compleja
interacción entre características físicas del inmueble, factores socioeconómicos
del entorno y, de manera fundamental, la \textbf{accesibilidad y proximidad a
servicios, infraestructura y equipamiento urbano}. Sin embargo, en Chile, y
particularmente en el Gran Santiago, los modelos de valoración inmobiliaria
disponibles públicamente suelen ignorar esta dimensión espacial, basándose
únicamente en atributos internos del predio o en tasaciones catastrales que
no capturan la dinámica real del mercado.

Este proyecto nace de la confluencia de dos disciplinas dentro del equipo:
\textbf{Planificación Urbana} y \textbf{Ciencia de Datos}. Uno de los integrantes
trabaja actualmente con geodatos urbanos en el marco de proyectos de planificación
territorial, mientras que los demás aportan experiencia en modelamiento estadístico
y aprendizaje de máquina. Esta combinación permite abordar la problemática desde
una perspectiva que integra el conocimiento del territorio con el rigor
cuantitativo del análisis de datos.

La motivación central es construir un modelo de predicción de precios inmobiliarios
que incorpore explícitamente la \textbf{dimensión geoespacial}: ¿cuánto vale un
predio según dónde está ubicado respecto a la red de metro, las paradas de
microbús, los centros educacionales, los establecimientos de salud, la actividad
comercial del entorno y otros factores urbanos medibles?

\section{Marco Teórico}

La valoración del suelo urbano se sustenta en un conjunto de teorías económicas
y espaciales que explican cómo la localización, la accesibilidad y la estructura
urbana determinan el precio de los predios. A continuación se presentan los
principales marcos conceptuales que fundamentan este proyecto.

\subsection{El Suelo como Bien Único}

A diferencia de otros bienes transables, el suelo posee una naturaleza
intrínsecamente singular que lo aleja del concepto de \textit{commodity}. Sus
atributos se dividen en dos categorías fundamentales:

\begin{itemize}[leftmargin=1.5em, itemsep=0.3em]
  \item \textbf{Atributos físicos:} determinados por la naturaleza; incluyen
        topografía, composición del suelo, orientación y condiciones
        medioambientales. Son fijos e inamovibles.
  \item \textbf{Atributos de uso:} determinados por la sociedad a través
        de normativas, planes reguladores y decisiones de zonificación. Definen
        qué puede construirse, a qué densidad y con qué destino.
\end{itemize}

Dado que no existen dos predios exactamente iguales, pues cada uno ocupa una
posición irrepetible en el espacio, el suelo constituye un bien
\textbf{heterogéneo por definición}. La normativa urbanística amplifica esta
heterogeneidad al asignar distintos derechos de edificación a predios contiguos,
generando diferencias de valor que no responden únicamente a características
físicas, sino a decisiones regulatorias. Esta singularidad justifica el uso de
modelos geoespaciales que capturen la variación local del precio, en lugar de
modelos globales que asuman homogeneidad en el mercado.

\subsection{Teorías de la Renta del Suelo}

\subsubsection*{Teoría Centro-Periferia: Johann Heinrich von Thünen (1783--1850)}

Von Thünen formuló el primer modelo formal de renta del suelo, originalmente
aplicado al uso agrícola. Su principio central establece que el valor de un
predio decrece con la distancia al centro urbano (\textit{CBD}), en función de
los costos de transporte necesarios para acceder a ese centro. La \textbf{renta
de localización} resulta del diferencial entre el precio obtenido en el mercado
central y los costos de trasladar la producción o la actividad desde el predio
hasta dicho centro.

Aunque concebida para contextos rurales, esta teoría sienta las bases del
pensamiento moderno sobre gradientes de precios urbanos: a mayor accesibilidad
al centro, mayor valor del suelo. Este principio es directamente observable en
el Gran Santiago, donde los precios por metro cuadrado disminuyen de forma
generalizada hacia la periferia.

\subsubsection*{Teoría de la Maximización de la Renta: William Alonso (1964)}

Alonso extendió el modelo de von Thünen al contexto urbano mediante su
concepto de \textbf{\textit{bid-rent}} o curva de oferta de renta. En este
marco, distintos agentes económicos —hogares, comercios, industrias— compiten
por la localización urbana ofreciendo rentas que reflejan la utilidad que
obtienen de cada emplazamiento. El equilibrio del mercado determina qué uso
ocupa cada zona:

\begin{itemize}[leftmargin=1.5em, itemsep=0.3em]
  \item Las actividades comerciales, que obtienen mayor beneficio de la
        centralidad, pagan rentas más altas y se concentran en el centro.
  \item Los hogares de mayores ingresos pueden elegir entre centralidad y
        espacio, generando patrones mixtos de localización.
  \item La industria, con menor sensibilidad a la accesibilidad comercial,
        se desplaza hacia zonas periféricas donde la renta es menor.
\end{itemize}

El modelo de Alonso predice que el precio del suelo es una función decreciente
y convexa de la distancia al centro, lo que constituye la hipótesis central que
este proyecto busca contrastar y enriquecer con variables geoespaciales
adicionales.

\subsection{Economías de Aglomeración}

Las \textbf{economías de aglomeración} describen los beneficios de productividad
y bienestar que se generan cuando actividades económicas y personas se concentran
en el espacio. En el contexto inmobiliario, se manifiestan de dos formas
principales:

\begin{itemize}[leftmargin=1.5em, itemsep=0.3em]
  \item \textbf{Economías de urbanización:} beneficios derivados del tamaño
        total de la ciudad, como mayor oferta de trabajo, servicios especializados
        e infraestructura compartida.
  \item \textbf{Economías de localización:} beneficios para firmas o actividades
        del mismo sector al estar próximas entre sí, como acceso a proveedores,
        trabajadores especializados y derrames de conocimiento.
\end{itemize}

Desde el punto de vista del suelo, la aglomeración de comercios, servicios y
empleos en un área incrementa su atractivo y, por tanto, la disposición a pagar
por los predios adyacentes. La variable de \textit{conteo de patentes
comerciales} en radios de proximidad, incluida en el diseño de \textit{features}
de este proyecto, operacionaliza directamente este concepto.

\section{Publico}

\subsection{Audiencia Objetivo}

El análisis está dirigido a múltiples audiencias:

\begin{itemize}[leftmargin=1.5em, itemsep=0.3em]
  \item \textbf{Ciudadanos e instituciones públicas:} que necesitan comprender
        la valorización diferenciada del suelo para evaluar políticas de vivienda
        asequible, plusvalías asociadas a inversión pública y equidad territorial.
  \item \textbf{Agentes inmobiliarios:} que buscan herramientas de
        valoración más precisas para poder captura la plusvalía.
  \item \textbf{Investigadores y planificadores urbanos:} interesados en cuantificar
        el impacto de la infraestructura sobre el valor del suelo.
  \item \textbf{Inversionistas y desarrolladores:} que requieren estimaciones de
        precio para evaluar la viabilidad de proyectos.
\end{itemize}

% ════════════════════════════════════════════════════════════════════════════
\section{Objetivos}

\subsection{Objetivo General}

Desarrollar un modelo predictivo del precio de viviendas y predios urbanos en
el área metropolitana de Santiago, incorporando variables geoespaciales derivadas
de capas de infraestructura, equipamiento y usos de suelo.

\subsection{Objetivos Específicos}

\begin{enumerate}[leftmargin=1.5em, itemsep=0.4em]
  \item \textbf{Caracterizar espacialmente} el mercado inmobiliario del área
        de estudio, identificando patrones de valorización por zonas y relación
        con el equipamiento urbano.
  \item \textbf{Construir y evaluar modelos de regresión} que relacionen el precio
        del inmueble con variables de accesibilidad (distancia a metro, paraderos,
        colegios, etc.) y características estructurales del predio.
  \item \textbf{Cuantificar el efecto marginal} de cada tipo de equipamiento
        sobre el precio, respondiendo preguntas como: ¿Cuánto aumenta el precio
        al estar a 100 m de una estación de metro? ¿Afecta negativamente la
        cercanía a zonas de alta accidentabilidad laboral?
  \item \textbf{Diseñar una herramienta de consulta} que permita estimar el
        precio esperado de un predio según su ubicación y características.
\end{enumerate}

\subsection{Preguntas de Investigación}

\begin{itemize}[leftmargin=1.5em, itemsep=0.3em]
  \item ¿Qué variables geoespaciales son los mejores predictores del precio de
        un inmueble en el Gran Santiago?
  \item ¿Existe un efecto de ``prima de accesibilidad'' asociado a la cercanía
        al transporte público?
  \item ¿Cómo varía espacialmente la importancia de los distintos factores a lo
        largo de la ciudad (efectos heterogéneos)?
\end{itemize}

% ════════════════════════════════════════════════════════════════════════════
\section{Datos}

\subsection{Fuentes y Descripción General}

El proyecto se apoya en un conjunto de \textbf{geodatabases} de acceso público,
principalmente provenientes del \textbf{Instituto Nacional de Estadísticas (INE)},
la \textbf{Infraestructura de Datos Espaciales de Chile (IDE Chile)}, portales
de datos abiertos y fuentes complementarias de mercado. Todos los datos se
encuentran en formatos geoespaciales estándar (\texttt{.shp}, \texttt{.geojson},
\texttt{.gpkg}) y pueden ser procesados con herramientas de SIG de código abierto
como QGIS, GeoPandas o PostGIS.

\vspace{0.5em}
\begin{center}
\begin{tabular}{>{\bfseries}p{4cm} p{4.5cm} p{4.5cm}}
\toprule
\textbf{Capa} & \textbf{Descripción} & \textbf{Fuente} \\
\midrule
Precios inmobiliarios & Transacciones de compraventa de viviendas y predios & Portal Observatorio Mercado del Suelo / Scraping portales inmobiliarios \\
\addlinespace
Patentes comerciales & Localización georreferenciada de comercios habilitados & Municipios / IDE Chile \\
\addlinespace
Red de metro & Estaciones y líneas del Metro de Santiago & Metro S.A. / IDE Chile \\
\addlinespace
Transporte público & Recorridos y paraderos de RED (ex-Transantiago) & DTPM / IDE Chile \\
\addlinespace
Centros de salud & Hospitales, CESFAM, clínicas georref. & MINSAL / IDE Chile \\
\addlinespace
Educación & Colegios, liceos y universidades & Ministerio de Educación / IDE Chile \\
\addlinespace
Accidentabilidad laboral & Accidentes del trabajo georreferenciados & ACHS (datos parcialmente desactualizados) \\
\addlinespace
Datos censales & Información socioeconómica por manzana & INE Censo 2024 \\
\addlinespace
Redes Viales & Carreteras, vias expresas & BCN / IDE Chile \\
\addlinespace
Catastro de predios & Roles y atributos de predios urbanos & SII \\
\bottomrule
\end{tabular}
\end{center}

\subsection{Características de los Datos}

\begin{itemize}[leftmargin=1.5em, itemsep=0.3em]
  \item \textbf{Tipo:} Datos vectoriales de puntos, líneas y polígonos; tablas
        de atributos asociadas; datos tabulares del INE y SII.
  \item \textbf{Formato:} \texttt{.shp}, \texttt{.geojson}, \texttt{.gpkg},
        \texttt{.csv}, \texttt{.xlsx}.
  \item \textbf{Volumen:} Se estima trabajar con cientos de miles de registros
        de predios para el área metropolitana de Santiago, más las capas de
        infraestructura correspondientes.
  \item \textbf{Sistema de Referencia:} Se usará \textbf{WGS84 / UTM Zona 19S}
        (EPSG:32719) para garantizar consistencia métrica en los cálculos de
        distancia.
\end{itemize}

\subsection{Variable Objetivo}

El precio de venta por metro cuadrado (\$/m²) o el precio total de la transacción,
obtenido de portales de publicación inmobiliaria mediante \textit{web scraping}
(Portal Inmobiliario, Yapo, TocToc) o de registros del Conservador de Bienes
Raíces cuando estén disponibles.

% ════════════════════════════════════════════════════════════════════════════
\section{Diseño Tentativo}

\subsection{Ingeniería de Variables Geoespaciales}

A partir de las capas descritas se calcularán features de accesibilidad para
cada predio de la muestra:

\begin{itemize}[leftmargin=1.5em, itemsep=0.3em]
  \item \textbf{Distancia euclidiana, distancia Manhattan y por red vial} a la estación de metro,
        paradero de bus y establecimiento de salud/educación más cercano.
  \item \textbf{Conteo de amenidades} en radios de 250 m, 500 m y 1.000 m
        (patentes comerciales, equipamientos, etc.).
  \item \textbf{Índices de densidad} de accidentes laborales por zona.
  \item \textbf{Variables contextuales censales} agregadas a nivel de manzana
        o zona censal (ingreso promedio, hacinamiento, etc.).
\end{itemize}

\subsection{Modelos a Explorar}

El análisis seguirá una progresión metodológica desde modelos interpretables
hacia modelos de mayor capacidad predictiva:

\begin{enumerate}[leftmargin=1.5em, itemsep=0.4em]
  \item \textbf{Regresión Hedónica (OLS):} modelo de referencia que relaciona
        el precio con características físicas y de localización, siguiendo la
        tradición de Rosen (1974). Permite interpretar directamente el efecto
        marginal de cada variable.
  \item \textbf{Regresión Geográficamente Ponderada (GWR):} extensión espacial
        del OLS que permite que los coeficientes varíen en el espacio, capturando
        heterogeneidades locales.
  \item \textbf{Modelos de ensamble:} Random Forest y Gradient Boosting
        (XGBoost/LightGBM) para capturar relaciones no lineales e interacciones
        entre variables, con análisis de importancia de features vía SHAP.
  \item \textbf{Redes neuronales con embedding espacial} (exploración opcional):
        incorporar coordenadas o embeddings de vecindario para capturar dependencia
        espacial de manera implícita.
\end{enumerate}

\subsection{Evaluación y Validación}

Los modelos se evaluarán mediante \textbf{validación cruzada espacial}
(\textit{spatial k-fold}), evitando fuga de información geográfica entre
conjuntos de entrenamiento y prueba. Las métricas principales serán:
$R^2$, \textit{Root Mean Squared Error} (RMSE) y \textit{Mean Absolute
Percentage Error} (MAPE).

\subsection{Herramientas Tecnológicas}

\begin{itemize}[leftmargin=1.5em, itemsep=0.2em]
  \item \textbf{Lenguaje:} Python 3.10+
  \item \textbf{SIG y geodatos:} GeoPandas, Shapely, Fiona, PostGIS
  \item \textbf{Modelamiento:} scikit-learn, XGBoost, LightGBM, mgwr
  \item \textbf{Visualización:} Matplotlib, Seaborn, Kepler.gl / Folium
  \item \textbf{Entorno:} Jupyter Notebooks + GitHub para control de versiones
\end{itemize}

\section{Uso de LLM}

Para explorar y determinar qué opciones de modelos existen y son adecuados para la tarea propuesta,
se utilizó \textbf{Claude Sonnet 4.6} como asistente de modelamiento. A continuación se presenta la
conversación sostenida con el modelo.

\bigskip

% ---------------------------------------------------------------
% Paquetes requeridos (agregar en el preámbulo si no están):
%   \usepackage{mdframed, xcolor, fontawesome5, geometry, lmodern}
%   \usepackage[most]{tcolorbox}
%   \usepackage{enumitem}
%   \usepackage{parskip}
% ---------------------------------------------------------------

% ---------- Definición de colores ----------
\definecolor{userbg}{RGB}{241,240,235}
\definecolor{userborder}{RGB}{180,178,169}
\definecolor{aibg}{RGB}{230,241,251}
\definecolor{aiborder}{RGB}{55,138,221}
\definecolor{codeboxbg}{RGB}{245,244,240}
\definecolor{codeborder}{RGB}{200,198,190}
\definecolor{labelai}{RGB}{24,95,165}
\definecolor{labeluser}{RGB}{80,78,72}
\definecolor{kwcolor}{RGB}{127,119,221}   % purple keywords
\definecolor{strcolor}{RGB}{13,158,117}   % teal strings
\definecolor{cmtcolor}{RGB}{136,135,128}  % gray comments

% ---------- Entornos de burbuja ----------

% Burbuja del usuario (centrada, ancho completo, gris claro)
\newenvironment{usermsg}{%
  \begin{tcolorbox}[
    enhanced,
    colback=userbg,
    colframe=userborder,
    arc=8pt,
    boxrule=0.4pt,
    left=10pt, right=10pt, top=7pt, bottom=7pt,
    fontupper=\small,
    before upper={\textcolor{labeluser}{\textbf{\small Usuario}}\par\smallskip},
    before={\smallskip},
    after={\smallskip}
  ]
}{%
  \end{tcolorbox}
}

% Burbuja del asistente (izquierda, azul claro)
\newenvironment{aimsg}{%
  \begin{tcolorbox}[
    enhanced,
    colback=aibg,
    colframe=aiborder,
    arc=8pt,
    boxrule=0.4pt,
    left=10pt, right=10pt, top=7pt, bottom=7pt,
    fontupper=\small,
    before upper={\textcolor{labelai}{\textbf{\small Claude Sonnet 4.6}}\par\smallskip},
    before={\begin{flushleft}\begin{minipage}{0.88\linewidth}},
    after={\end{minipage}\end{flushleft}\smallskip}
  ]
}{%
  \end{tcolorbox}
}

% Bloque de código inline dentro de burbuja
\newcommand{\inlinecode}[1]{%
  \colorbox{codeboxbg}{\texttt{\small #1}}%
}

% Bloque de código multilínea dentro de burbuja
\newenvironment{codeblock}{%
  \begin{tcolorbox}[
    colback=codeboxbg,
    colframe=codeborder,
    arc=4pt,
    boxrule=0.3pt,
    left=8pt, right=8pt, top=5pt, bottom=5pt,
    fontupper=\ttfamily\footnotesize,
    before={\smallskip},
    after={\smallskip}
  ]
}{%
  \end{tcolorbox}
}

% ================================================================
%  CONVERSACIÓN
% ================================================================

\begin{usermsg}
Tenemos diversos shp y geopackages de líneas y paradas de metro, bus, centros de salud,
educacionales, comerciales, entre otros. ¿Qué modelos de modelamiento matemático--computacional
nos recomiendas para predecir los precios de diversos predios, asumiendo que el suelo es un bien
heterogéneo por definición?
\end{usermsg}

\begin{aimsg}
Dado que disponen de capas SHP y GeoPackages de transporte, equipamiento y comercio, el camino
metodológico natural se organiza en cuatro familias de modelos, siguiendo una progresión desde
modelos interpretables hacia modelos de mayor capacidad predictiva:

\medskip
\textbf{Modelos a explorar:}

\begin{enumerate}[leftmargin=1.5em, itemsep=0.4em]
  \item \textbf{Regresión Hedónica (OLS):} modelo de referencia que relaciona
        el precio con características físicas y de localización, siguiendo la
        tradición de Rosen (1974). Permite interpretar directamente el efecto
        marginal de cada variable.
  \item \textbf{Regresión Geográficamente Ponderada (GWR):} extensión espacial
        del OLS que permite que los coeficientes varíen en el espacio, capturando
        heterogeneidades locales.
  \item \textbf{Modelos de ensamble:} Random Forest y Gradient Boosting
        (XGBoost/LightGBM) para capturar relaciones no lineales e interacciones
        entre variables, con análisis de importancia de features vía SHAP.
  \item \textbf{Redes neuronales con embedding espacial} (exploración opcional):
        incorporar coordenadas o embeddings de vecindario para capturar dependencia
        espacial de manera implícita.
\end{enumerate}

\medskip
\textbf{Evaluación y validación:}

Los modelos se evaluarán mediante \textbf{validación cruzada espacial}
(\textit{spatial k-fold}), evitando fuga de información geográfica entre
conjuntos de entrenamiento y prueba. Las métricas principales serán:
$R^2$, \textit{Root Mean Squared Error} (RMSE) y \textit{Mean Absolute
Percentage Error} (MAPE).

\medskip
\textbf{Herramientas tecnológicas:}

\begin{itemize}[leftmargin=1.5em, itemsep=0.2em]
  \item \textbf{Lenguaje:} Python 3.10+
  \item \textbf{SIG y geodatos:} GeoPandas, Shapely, Fiona, PostGIS
  \item \textbf{Modelamiento:} scikit-learn, XGBoost, LightGBM, mgwr
  \item \textbf{Visualización:} Matplotlib, Seaborn, Kepler.gl / Folium
  \item \textbf{Entorno:} Jupyter Notebooks + GitHub para control de versiones
\end{itemize}
\end{aimsg}

\begin{usermsg}
¿Cuáles son las variables de entrada más relevantes que podemos extraer directamente desde
nuestros archivos SHP y GeoPackages?
\end{usermsg}

\begin{aimsg}
Desde las capas disponibles es posible derivar cuatro grupos de variables de localización:

\medskip
\textbf{Accesibilidad a transporte:}
\begin{itemize}
  \item Distancia euclidiana y por red vial a la parada de metro o bus más cercana.
  \item Número de líneas de bus en un radio de 500\,m y 1\,000\,m (\textit{buffer count}).
  \item Índice de frecuencia de servicio ponderado por distancia.
\end{itemize}

\textbf{Equipamiento de salud y educación:}
\begin{itemize}
  \item Distancia al centro de salud primario más cercano (CESFAM, consultorio).
  \item Número de establecimientos educacionales (básica, media) en radio de 1\,km.
  \item Clasificación del establecimiento más cercano (público / subvencionado / particular).
\end{itemize}

\textbf{Centralidad comercial:}
\begin{itemize}
  \item Densidad de comercio en radio de 500\,m (puntos por $\text{km}^2$).
  \item Distancia al nodo comercial de mayor densidad (centro de barrio).
\end{itemize}

\textbf{Entorno físico y uso de suelo:}
\begin{itemize}
  \item Categoría de uso de suelo dominante en buffer de 250\,m.
  \item Superficie de áreas verdes accesibles en radio de 800\,m.
  \item Pendiente media del terreno (si se dispone de DEM).
\end{itemize}

\medskip
Todas estas variables se calculan con \inlinecode{geopandas} mediante operaciones de
\inlinecode{sjoin}, \inlinecode{buffer} y \inlinecode{distance}, o equivalentemente en
\inlinecode{PostGIS} con funciones \inlinecode{ST\_DWithin} y \inlinecode{ST\_Distance}.
\end{aimsg}

\end{document}
